%% Section 2: Background and Related Work
%% English manuscript section.
%% -----------------------------------------------------------------------

\section{Background and Related Work}
\label{sec:background}

\subsection{AI for Mathematics}
\label{sec:background:ai-math}

Throughout mathematics, numerical pattern observation has seeded important
conjectures---from Monstrous Moonshine~\cite{conway1979monstrous} and the
BSD conjecture~\cite{birch1965} to Wiles's proof of the Shimura--Taniyama conjecture
for semistable elliptic curves, which resolved Fermat's Last Theorem~\cite{wiles1995},
and Candelas et al.'s mirror-symmetry prediction of 317,206,375
rational curves~\cite{candelas1991}.
Recent AI systems including AlphaGeometry~\cite{trinh2024alphageometry},
AlphaProof~\cite{alphaproof2024}, the Ramanujan Machine~\cite{raayoni2021ramanujan}, and
FunSearch~\cite{romera2023funsearch} demonstrate the viability of AI-driven mathematical
discovery. On the formalisation front, Urban's large-scale autoformalization
work~\cite{urban-autoformalization} shows that natural-language mathematics can be
automatically translated into machine-verifiable form.
Yet extracting arithmetic laws from integer sequences at scale remains largely
unexplored. This work aims to establish representational foundations for that capability.

\subsection{AI Research on OEIS}
\label{sec:background:oeis}

Program-synthesis approaches to OEIS include Alien Coding~\cite{urban-alien-coding},
which searches for the shortest reproducing code; QSynt~\cite{gauthier2023qsynt}, which
synthesises programs for 43,516 sequences via self-learning tree search; and Learning
Conjecturing~\cite{gauthier2025conjecturing}, which automatically proves 5,565 arithmetic problems.
FACT~\cite{zurich-fact} (Bel\v{c}\'{a}k et al., NeurIPS 2022) constructs a comprehensive OEIS
benchmark covering five tasks in order of difficulty---classification, similarity,
next sequence-part prediction, continuation, and
\emph{unmasking} (identified as the hardest).
Our Vanilla baseline corresponds to FACT's plain Transformer; IntSeqBERT extends it with
modulo feature engineering and FiLM fusion to overcome FACT's limitations with large integers
and multiplicative structure. The Solver-based next-term prediction evaluated in
Section~\ref{sec:solver-eval} corresponds to FACT's continuation task.

\subsection{Modulo Spectra as Arithmetic Feature Engineering}
\label{sec:background:modulo-engineering}

Deep networks exhibit a well-known \emph{spectral bias}~\cite{rahaman2019spectral}:
they learn low-frequency functions before high-frequency ones. In the context of integer
sequences, this is particularly problematic for multiplicatively growing sequences
(e.g., $n^2$, $n!$), where learning the growth law from token sequences alone requires
many layers.
This work mitigates the problem through feature engineering: because the results of
multiplication appear naturally in modular residues, supplying the modulo spectrum as explicit
input features reduces the network depth required to ``rediscover'' multiplicative structure.

More precisely, most OEIS sequences are generated by finite algorithms combining addition,
multiplication, and modular operations. Since modular arithmetic is a ring homomorphism over
$\mathbb{Z}$, the residue sequence $(x_i \bmod m)$ faithfully preserves additive and
multiplicative structure as a compact representation, regardless of the magnitudes involved.
Fermat's little theorem implies that $(a^n \bmod p)$ is periodic with period dividing $p-1$
for prime $p$ with $\gcd(a,p)=1$. Gauss's law of quadratic
reciprocity~\cite{gauss-reciprocity} and its generalisations further link residue information
across different primes.
Composite moduli simultaneously retain information from all prime power factors via CRT,
motivating the use of both prime and composite moduli in our spectrum.
Our modulo features cover $m \in \{2, 3, \ldots, 101\}$ (100 moduli), capturing power-residue
structure broadly across primes and composites.

\subsection{Positioning and Neural Representations}
\label{sec:background:repr-position}

Tokenisation~(the LLM mainstream) cannot handle out-of-vocabulary values and buries
arithmetic structure in opaque token IDs. Log-scale continuous embeddings improve scale
invariance but miss arithmetic periodicity. We instead repurpose sinusoidal
embeddings~\cite{vaswani2017attention} from positions to \emph{modular residues of values},
and apply FiLM~\cite{perez2018film} to modulate the magnitude stream via the modulo stream.


In contrast to the Gauthier and Urban group's program-synthesis line~\cite{urban-alien-coding,gauthier2023qsynt,gauthier2025conjecturing},
which explicitly describes each sequence's generation rule, this work uses \emph{representation
learning} to acquire shared arithmetic structure across the OEIS corpus.
We view the two approaches as complementary: program synthesis seeks a formula or program
that generates a given sequence, whereas our representations aim to judge whether two
sequences are number-theoretically similar even when they are generated in entirely
different ways. Pre-trained IntSeqBERT embeddings could in turn serve as a search heuristic
for program synthesis, or as an index for retrieving similar sequences.
Where FACT~\cite{zurich-fact} identified the limits of a plain tokenised Transformer,
IntSeqBERT overcomes them via modulo-spectrum feature engineering and FiLM fusion.
